\section{Photonic Formulation}

In photonic quantum computing, we use $n$ modes to represent $n$ vertices. We inject exactly $k = n/2$ photons into the first $k$ modes.

A passive linear interferometer $\hat{U}$ conserves the total photon number:

\begin{align}
    [\hat{U}, \hat{N}] = 0
\end{align}

where $\hat{N} = \sum_i a_i^\dagger a_i$ is the photon number operator.

Therefore, if the input state has $k$ photons, the output state also has $k$ photons:

\begin{align}
    \sum_i n_i^{\text{out}} = k
\end{align}

\textbf{Critical distinction:} Fixed photon number does NOT imply collision-free occupations.

The full $k$-photon Fock sector is:
\begin{align}
    \mathcal{H}_k = \text{span}\{|n_1,\ldots,n_n\rangle : \sum_i n_i = k\}
\end{align}

The collision-free subspace is:
\begin{align}
    \mathcal{H}_k^{CF} = \text{span}\{|n_1,\ldots,n_n\rangle : \sum_i n_i = k, n_i \in \{0,1\}\}
\end{align}

For $m$ modes and $k$ photons:
\begin{align}
    \dim(\mathcal{H}_k) &= \binom{m+k-1}{k} \\
    \dim(\mathcal{H}_k^{CF}) &= \binom{m}{k}
\end{align}

\textbf{Key insight:} Collision-free outcomes map exactly to balanced partitions, but $P_{CF} < 1$ for generic interferometers.
